\documentclass[12pt, letterpaper]{article}
\usepackage[T1]{fontenc}
\usepackage[spanish]{babel}
\usepackage[utf8]{inputenc}
\usepackage{graphicx}
\usepackage{geometry}
\usepackage{titlesec}
\usepackage{fancyhdr}
\usepackage{setspace}
\usepackage{enumitem}
\usepackage{tocloft}
\usepackage{times}
\usepackage{float}
\usepackage{hyperref}
\usepackage{listings}
\usepackage{xcolor}

\definecolor{codegreen}{rgb}{0.12, 0.53, 0.23}
\definecolor{codeblue}{rgb}{0.01, 0.40, 0.84}
\definecolor{codepurple}{rgb}{0.50, 0.00, 0.50}
\definecolor{codebackground}{rgb}{0.97, 0.98, 0.99}
\definecolor{codegray}{rgb}{0.45, 0.45, 0.45}
\definecolor{codeorange}{rgb}{0.85, 0.45, 0.0}

\lstdefinelanguage{purePHP}{
    sensitive=true,
    keywords={class,extends,public,protected,private,function,return,use,namespace,if,else,foreach,endforeach,endif,or,and,as,array,static,new,try,catch,throw,compact,route,view,dd,old,session,schema,blueprint,extends,implements,interface,trait,string,int,float,bool,void},
    morecomment=[l]{//},
    morecomment=[s]{/*}{*/},
    morestring=[b]",
    morestring=[b]',
}

\hypersetup{
    hypertexnames=false,
    colorlinks=true,
    linkcolor=black,
    urlcolor=blue
}

% Configuración básica
\geometry{letterpaper, left=3cm, right=2.5cm, top=2.5cm, bottom=2.5cm}
\renewcommand{\baselinestretch}{1.5}
\setlength{\parindent}{0pt}
\setlength{\parskip}{1em}

% Configuración de numeración
\fancypagestyle{plain}{\fancyhf{}\fancyfoot[R]{\thepage}}
\pagestyle{plain}

\setcounter{secnumdepth}{4}

% Configuración de títulos
\titleformat{\section}{\raggedleft\bfseries}{}{0pt}{}
\titleformat{\subsection}{\centering\bfseries}{\arabic{subsection}.}{0.5em}{}
\titleformat{\subsubsection}{\raggedright\bfseries}{\arabic{subsection}.\arabic{subsubsection}.}{0.5em}{}

\renewcommand{\thesubsection}{\arabic{subsection}}
\renewcommand{\thesubsubsection}{\arabic{subsection}.\arabic{subsubsection}}

\renewcommand{\cftsecfont}{\bfseries}
\renewcommand{\cftsubsecfont}{\bfseries}
\renewcommand{\cftsubsubsecfont}{\bfseries}

\fontsize{14}{18}\selectfont

% Configuración de Listings para código limpio en español
\lstset{
    basicstyle=\ttfamily\footnotesize,
    backgroundcolor=\color{codebackground},
    commentstyle=\color{codegreen}\itshape,
    keywordstyle=\color{codeblue}\bfseries,
    numberstyle=\tiny\color{codegray},
    stringstyle=\color{codepurple},
    breaklines=true,
    showstringspaces=false,
    numbers=left,
    frame=single,
    rulecolor=\color{gray!30},
    extendedchars=true,
    inputencoding=utf8,
    literate={á}{{\'a}}1 {é}{{\'e}}1 {í}{{\'i}}1 {ó}{{\'o}}1 {ú}{{\'u}}1 {ñ}{{\~n}}1 {Á}{{\'A}}1 {É}{{\'E}}1 {Í}{{\'I}}1 {Ó}{{\'O}}1 {Ú}{{\'U}}1 {Ñ}{{\~N}}1 {¿}{{\textquestiondown}}1 {¡}{{\textexclamdown}}1 {🍗}{{\raisebox{-0.25ex}{\includegraphics[height=2.2ex]{chicken.png}}}}1
}

\begin{document}

% Portada
\title{
    \vspace*{-2cm}
    UNIVERSIDAD PRIVADA DOMINGO SAVIO\\
    SEDE SUCRE\\
    \vspace{0.5cm}
    \includegraphics[width=0.4\textwidth]{logo.png}\\
    \vspace{0.5cm} 
    INGENIERÍA DE SISTEMAS\\
    \vspace{0.75cm}
    Laboratorio 4: Segunda presentación Sistema de venta broasteria
}
\author{
    Pedro Antonio López Chumacero
}
\date{
    MATERIA: Desarrollo de Sistemas \\
    DOCENTE: Ing. Erick Vidal Juarez Choque \\
    \vspace{0.75cm}
    SUCRE - BOLIVIA \\
    \today
}

\begin{singlespace}
\maketitle
\thispagestyle{empty}
\end{singlespace}
\newpage
\clearpage

\pagenumbering{roman}
\tableofcontents
\thispagestyle{empty}
\newpage
\clearpage

\listoffigures
\thispagestyle{empty}
\newpage
\clearpage

\pagenumbering{arabic}
\setcounter{page}{1}

\section*{CONFIGURACIÓN GENERAL Y RUTAS}
\addcontentsline{toc}{section}{CONFIGURACIÓN GENERAL Y RUTAS}
\setcounter{section}{1}

\subsection{Rutas del Sistema (routes/web.php)}
\begin{lstlisting}[language=purePHP]
<?php

use Illuminate\Support\Facades\Route;

use App\Http\Controllers\TurnoController;
use App\Http\Controllers\EmpleadoController;
use App\Http\Controllers\ClienteController;
use App\Http\Controllers\InventarioController;
use App\Http\Controllers\PedidoController;
use App\Http\Controllers\DetallePedidoController;
use App\Http\Controllers\VentaController;

Route::get('/', function () {
    return redirect()->route('empleados.index');
});

Route::resource('turnos', TurnoController::class);
Route::resource('empleados', EmpleadoController::class);
Route::resource('clientes', ClienteController::class);
Route::resource('inventarios', InventarioController::class);
Route::resource('pedidos', PedidoController::class);
Route::resource('detalle-pedidos', DetallePedidoController::class);
Route::resource('ventas', VentaController::class);
\end{lstlisting}

\subsection{Diseño Base (resources/views/layouts/app.blade.php)}
\begin{lstlisting}[language=HTML]
<!DOCTYPE html>
<html lang="es">
<head>
    <meta charset="UTF-8">
    <meta name="viewport" content="width=device-width, initial-scale=1.0">
    <title>@yield('title', 'Pollos Luisa') - Panel de Gestion</title>
    <link rel="preconnect" href="https://fonts.googleapis.com">
    <link rel="preconnect" href="https://fonts.gstatic.com" crossorigin>
    <link href="https://fonts.googleapis.com/css2?family=Inter:wght@300;400;500;600;700&display=swap" rel="stylesheet">
    <style>
        :root {
            --bg-color: #0b0f19;
            --card-bg: rgba(255, 255, 255, 0.03);
            --card-border: rgba(255, 255, 255, 0.08);
            --primary-color: #ff5e3a;
            --primary-hover: #ff451b;
            --secondary-color: #4f46e5;
            --secondary-hover: #4338ca;
            --text-main: #f3f4f6;
            --text-muted: #9ca3af;
            --accent-green: #10b981;
            --accent-red: #ef4444;
            --glass-blur: blur(12px);
        }

        * { box-sizing: border-box; margin: 0; padding: 0; }

        body {
            font-family: 'Inter', sans-serif;
            background-color: var(--bg-color);
            background-image: 
                radial-gradient(at 0% 0%, rgba(79, 70, 229, 0.15) 0px, transparent 50%),
                radial-gradient(at 100% 100%, rgba(255, 94, 58, 0.1) 0px, transparent 50%);
            color: var(--text-main);
            min-height: 100vh;
            display: flex;
            flex-direction: column;
            line-height: 1.5;
        }

        a { color: inherit; text-decoration: none; }

        header {
            background: rgba(11, 15, 25, 0.7);
            backdrop-filter: var(--glass-blur);
            border-bottom: 1px solid var(--card-border);
            position: sticky; top: 0; z-index: 100;
        }

        .nav-container {
            max-width: 1200px; margin: 0 auto; padding: 1rem 2rem;
            display: flex; justify-content: space-between; align-items: center;
        }

        .logo {
            font-size: 1.5rem; font-weight: 700;
            background: linear-gradient(135deg, var(--primary-color), #ff9e80);
            -webkit-background-clip: text; -webkit-text-fill-color: transparent;
            display: flex; align-items: center; gap: 0.5rem;
        }

        .nav-links { display: flex; gap: 1.5rem; }

        .nav-link {
            font-size: 0.95rem; font-weight: 500; color: var(--text-muted);
            padding: 0.5rem 1rem; border-radius: 8px; transition: all 0.3s ease;
        }

        .nav-link:hover, .nav-link.active {
            color: var(--text-main); background: rgba(255, 255, 255, 0.05);
        }

        main { flex: 1; max-width: 1200px; width: 100%; margin: 2rem auto; padding: 0 2rem; }

        .glass-card {
            background: var(--card-bg); backdrop-filter: var(--glass-blur);
            border: 1px solid var(--card-border); border-radius: 16px; padding: 2rem;
            box-shadow: 0 10px 30px rgba(0, 0, 0, 0.5); margin-bottom: 2rem;
        }

        .card-header {
            display: flex; justify-content: space-between; align-items: center;
            margin-bottom: 2rem; border-bottom: 1px solid var(--card-border); padding-bottom: 1rem;
        }

        .card-title { font-size: 1.5rem; font-weight: 600; color: var(--text-main); }

        .alert {
            padding: 1rem; border-radius: 8px; margin-bottom: 1.5rem;
            display: flex; align-items: center; font-size: 0.95rem; border: 1px solid transparent;
        }

        .alert-success { background: rgba(16, 185, 129, 0.1); border-color: rgba(16, 185, 129, 0.2); color: #34d399; }
        .alert-error { background: rgba(239, 68, 68, 0.1); border-color: rgba(239, 68, 68, 0.2); color: #f87171; }

        .btn {
            display: inline-flex; align-items: center; justify-content: center;
            padding: 0.6rem 1.2rem; font-size: 0.9rem; font-weight: 600; border-radius: 8px;
            cursor: pointer; transition: all 0.2s ease; border: none; gap: 0.5rem;
        }

        .btn-primary { background: var(--primary-color); color: #ffffff; }
        .btn-primary:hover { background: var(--primary-hover); transform: translateY(-1px); }
        .btn-secondary { background: var(--secondary-color); color: #ffffff; }
        .btn-secondary:hover { background: var(--secondary-hover); transform: translateY(-1px); }
        .btn-outline { background: transparent; color: var(--text-main); border: 1px solid var(--card-border); }
        .btn-outline:hover { background: rgba(255, 255, 255, 0.05); border-color: var(--text-muted); }
        .btn-danger { background: rgba(239, 68, 68, 0.15); color: #f87171; border: 1px solid rgba(239, 68, 68, 0.25); }
        .btn-danger:hover { background: var(--accent-red); color: #ffffff; }

        .table-responsive { overflow-x: auto; width: 100%; }
        table { width: 100%; border-collapse: collapse; text-align: left; }
        th { padding: 1rem; font-size: 0.85rem; font-weight: 600; text-transform: uppercase; letter-spacing: 0.05em; color: var(--text-muted); border-bottom: 1px solid var(--card-border); }
        td { padding: 1.2rem 1rem; font-size: 0.95rem; border-bottom: 1px solid var(--card-border); color: var(--text-main); }
        tr:hover td { background: rgba(255, 255, 255, 0.01); }
        .actions-cell { display: flex; gap: 0.5rem; }

        .form-group { margin-bottom: 1.5rem; }
        .form-label { display: block; margin-bottom: 0.5rem; font-size: 0.9rem; font-weight: 500; color: var(--text-main); }
        .form-control { width: 100%; padding: 0.75rem 1rem; font-size: 0.95rem; background: rgba(0, 0, 0, 0.2); border: 1px solid var(--card-border); border-radius: 8px; color: var(--text-main); outline: none; transition: all 0.2s ease; }
        .form-control:focus { border-color: var(--primary-color); box-shadow: 0 0 0 2px rgba(255, 94, 58, 0.15); }
        .error-message { color: #f87171; font-size: 0.8rem; margin-top: 0.25rem; display: block; }

        footer { text-align: center; padding: 2rem; border-top: 1px solid var(--card-border); color: var(--text-muted); font-size: 0.85rem; background: rgba(11, 15, 25, 0.5); }
        .badge { display: inline-flex; align-items: center; padding: 0.25rem 0.75rem; font-size: 0.75rem; font-weight: 600; border-radius: 50px; background: rgba(255, 255, 255, 0.05); color: var(--text-muted); }
        .badge-accent { background: rgba(79, 70, 229, 0.2); color: #a5b4fc; }
    </style>
</head>
<body>
    <header>
        <div class="nav-container">
            <a href="{{ url('/') }}" class="logo">Pollos Luisa</a>
            <nav class="nav-links">
                <a href="{{ route('empleados.index') }}" class="nav-link {{ Request::is('empleados*') ? 'active' : '' }}">Empleados</a>
                <a href="{{ route('turnos.index') }}" class="nav-link {{ Request::is('turnos*') ? 'active' : '' }}">Turnos</a>
                <a href="{{ route('clientes.index') }}" class="nav-link {{ Request::is('clientes*') ? 'active' : '' }}">Clientes</a>
                <a href="{{ route('inventarios.index') }}" class="nav-link {{ Request::is('inventarios*') ? 'active' : '' }}">Inventarios</a>
                <a href="{{ route('pedidos.index') }}" class="nav-link {{ Request::is('pedidos*') ? 'active' : '' }}">Pedidos</a>
                <a href="{{ route('detalle-pedidos.index') }}" class="nav-link {{ Request::is('detalle-pedidos*') ? 'active' : '' }}">Detalle Pedidos</a>
                <a href="{{ route('ventas.index') }}" class="nav-link {{ Request::is('ventas*') ? 'active' : '' }}">Ventas</a>
            </nav>
        </div>
    </header>
    <main>
        @if(session('success')) <div class="alert alert-success">{{ session('success') }}</div> @endif
        @if(session('error')) <div class="alert alert-error">{{ session('error') }}</div> @endif
        @yield('content')
    </main>
    <footer>&copy; {{ date('Y') }} Pollos Luisa. Todos los derechos reservados.</footer>
</body>
</html>
\end{lstlisting}

\newpage

\section*{CRUD CLIENTES}
\addcontentsline{toc}{section}{CRUD CLIENTES}
\setcounter{section}{2}

\subsection{Migración (2026\_07\_22\_000002\_create\_clientes\_table.php)}
\begin{lstlisting}[language=purePHP]
<?php

use Illuminate\Database\Migrations\Migration;
use Illuminate\Database\Schema\Blueprint;
use Illuminate\Support\Facades\Schema;

return new class extends Migration
{
    public function up(): void
    {
        Schema::create('clientes', function (Blueprint $table) {
            $table->id();
            $table->integer('ci');
            $table->string('nombre_completo');
            $table->string('email');
            $table->integer('telefono');
            $table->date('fecha_nacimiento');
            $table->timestamps();
        });
    }

    public function down(): void
    {
        Schema::dropIfExists('clientes');
    }
};
\end{lstlisting}

\subsection{Modelo (Cliente.php)}
\begin{lstlisting}[language=purePHP]
<?php

namespace App\Models;

use Illuminate\Database\Eloquent\Factories\HasFactory;
use Illuminate\Database\Eloquent\Model;

class Cliente extends Model
{
    use HasFactory;

    protected $table = 'clientes';

    protected $fillable = [
        'ci',
        'nombre_completo',
        'email',
        'telefono',
        'fecha_nacimiento',
    ];

    public function pedidos()
    {
        return $this->hasMany(Pedido::class, 'cliente_id');
    }
}
\end{lstlisting}

\subsection{Controlador (ClienteController.php)}
\begin{lstlisting}[language=purePHP]
<?php

namespace App\Http\Controllers;

use App\Models\Cliente;
use Illuminate\Http\Request;

class ClienteController extends Controller
{
    public function index()
    {
        $clientes = Cliente::latest()->get();
        return view('clientes.index', compact('clientes'));
    }

    public function create()
    {
        return view('clientes.create');
    }

    public function store(Request $request)
    {
        $validated = $request->validate([
            'ci' => 'required|integer',
            'nombre_completo' => 'required|string|max:255',
            'email' => 'required|email|max:255',
            'telefono' => 'required|integer',
            'fecha_nacimiento' => 'required|date',
        ]);

        Cliente::create($validated);

        return redirect()->route('clientes.index')->with('success', 'Cliente registrado con exito.');
    }

    public function show(Cliente $cliente)
    {
        $cliente->load('pedidos');
        return view('clientes.show', compact('cliente'));
    }

    public function edit(Cliente $cliente)
    {
        return view('clientes.edit', compact('cliente'));
    }

    public function update(Request $request, Cliente $cliente)
    {
        $validated = $request->validate([
            'ci' => 'required|integer',
            'nombre_completo' => 'required|string|max:255',
            'email' => 'required|email|max:255',
            'telefono' => 'required|integer',
            'fecha_nacimiento' => 'required|date',
        ]);

        $cliente->update($validated);

        return redirect()->route('clientes.index')->with('success', 'Cliente actualizado con exito.');
    }

    public function destroy(Cliente $cliente)
    {
        $cliente->delete();

        return redirect()->route('clientes.index')->with('success', 'Cliente eliminado con exito.');
    }
}
\end{lstlisting}

\subsection{Vista Principal (clientes/index.blade.php)}
\begin{lstlisting}[language=HTML]
@extends('layouts.app')

@section('title', 'Clientes')

@section('content')
<div class="glass-card">
    <div class="card-header">
        <div>
            <h1 class="card-title">Gestion de Clientes</h1>
            <p style="color: var(--text-muted); font-size: 0.9rem;">Registro y administracion de clientes</p>
        </div>
        <a href="{{ route('clientes.create') }}" class="btn btn-primary">
            + Nuevo Cliente
        </a>
    </div>

    <div class="table-responsive">
        <table>
            <thead>
                <tr>
                    <th>ID</th>
                    <th>CI</th>
                    <th>Nombre Completo</th>
                    <th>Email</th>
                    <th>Telefono</th>
                    <th>Fecha Nacimiento</th>
                    <th>Acciones</th>
                </tr>
            </thead>
            <tbody>
                @forelse($clientes as $cli)
                <tr>
                    <td>#{{ $cli->id }}</td>
                    <td>{{ $cli->ci }}</td>
                    <td><strong>{{ $cli->nombre_completo }}</strong></td>
                    <td>{{ $cli->email }}</td>
                    <td>{{ $cli->telefono }}</td>
                    <td>{{ \Carbon\Carbon::parse($cli->fecha_nacimiento)->format('d/m/Y') }}</td>
                    <td>
                        <div class="actions-cell">
                            <a href="{{ route('clientes.show', $cli->id) }}" class="btn btn-outline" style="padding: 0.3rem 0.6rem; font-size: 0.8rem;">Ver</a>
                            <a href="{{ route('clientes.edit', $cli->id) }}" class="btn btn-secondary" style="padding: 0.3rem 0.6rem; font-size: 0.8rem;">Editar</a>
                            <form action="{{ route('clientes.destroy', $cli->id) }}" method="POST" onsubmit="return confirm('Seguro de eliminar este cliente?');" style="display:inline;">
                                @csrf
                                @method('DELETE')
                                <button type="submit" class="btn btn-danger" style="padding: 0.3rem 0.6rem; font-size: 0.8rem;">Eliminar</button>
                            </form>
                        </div>
                    </td>
                </tr>
                @empty
                <tr>
                    <td colspan="7" style="text-align: center; color: var(--text-muted); padding: 2rem;">
                        No hay clientes registrados.
                    </td>
                </tr>
                @endforelse
            </tbody>
        </table>
    </div>
</div>
@endsection
\end{lstlisting}

\subsection{Vista del CRUD de Clientes en Funcionamiento}
A continuacion se presenta la captura de pantalla de la interfaz del CRUD de Clientes:

\begin{figure}[H]
    \centering
    \includegraphics[width=0.9\textwidth]{view-clientes.png}
    \caption{Interfaz de usuario del CRUD de Clientes}
\end{figure}

\newpage

\section*{CRUD INVENTARIOS}
\addcontentsline{toc}{section}{CRUD INVENTARIOS}
\setcounter{section}{3}

\subsection{Migración (2026\_07\_22\_000001\_create\_inventarios\_table.php)}
\begin{lstlisting}[language=purePHP]
<?php

use Illuminate\Database\Migrations\Migration;
use Illuminate\Database\Schema\Blueprint;
use Illuminate\Support\Facades\Schema;

return new class extends Migration
{
    public function up(): void
    {
        Schema::create('inventarios', function (Blueprint $table) {
            $table->id();
            $table->string('producto');
            $table->decimal('stock_actual', 10, 2)->default(0);
            $table->decimal('stock_minimo', 10, 2)->default(0);
            $table->decimal('precio', 10, 2)->default(0.00);
            $table->timestamps();
        });
    }

    public function down(): void
    {
        Schema::dropIfExists('inventarios');
    }
};
\end{lstlisting}

\subsection{Modelo (Inventario.php)}
\begin{lstlisting}[language=purePHP]
<?php

namespace App\Models;

use Illuminate\Database\Eloquent\Factories\HasFactory;
use Illuminate\Database\Eloquent\Model;

class Inventario extends Model
{
    use HasFactory;

    protected $table = 'inventarios';

    protected $fillable = [
        'producto',
        'stock_actual',
        'stock_minimo',
        'precio',
    ];

    public function detallePedidos()
    {
        return $this->hasMany(DetallePedido::class, 'inventario_id');
    }
}
\end{lstlisting}

\subsection{Controlador (InventarioController.php)}
\begin{lstlisting}[language=purePHP]
<?php

namespace App\Http\Controllers;

use App\Models\Inventario;
use Illuminate\Http\Request;

class InventarioController extends Controller
{
    public function index()
    {
        $inventarios = Inventario::latest()->get();
        return view('inventarios.index', compact('inventarios'));
    }

    public function create()
    {
        return view('inventarios.create');
    }

    public function store(Request $request)
    {
        $validated = $request->validate([
            'producto' => 'required|string|max:255',
            'stock_actual' => 'required|numeric|min:0',
            'stock_minimo' => 'required|numeric|min:0',
            'precio' => 'required|numeric|min:0',
        ]);

        Inventario::create($validated);

        return redirect()->route('inventarios.index')->with('success', 'Producto registrado en el inventario con exito.');
    }

    public function show(Inventario $inventario)
    {
        return view('inventarios.show', compact('inventario'));
    }

    public function edit(Inventario $inventario)
    {
        return view('inventarios.edit', compact('inventario'));
    }

    public function update(Request $request, Inventario $inventario)
    {
        $validated = $request->validate([
            'producto' => 'required|string|max:255',
            'stock_actual' => 'required|numeric|min:0',
            'stock_minimo' => 'required|numeric|min:0',
            'precio' => 'required|numeric|min:0',
        ]);

        $inventario->update($validated);

        return redirect()->route('inventarios.index')->with('success', 'Producto de inventario actualizado con exito.');
    }

    public function destroy(Inventario $inventario)
    {
        $inventario->delete();

        return redirect()->route('inventarios.index')->with('success', 'Producto eliminado del inventario.');
    }
}
\end{lstlisting}

\subsection{Vista Principal (inventarios/index.blade.php)}
\begin{lstlisting}[language=HTML]
@extends('layouts.app')

@section('title', 'Inventarios')

@section('content')
<div class="glass-card">
    <div class="card-header">
        <div>
            <h1 class="card-title">Gestion de Inventario</h1>
            <p style="color: var(--text-muted); font-size: 0.9rem;">Control de insumos, productos y stock disponible</p>
        </div>
        <a href="{{ route('inventarios.create') }}" class="btn btn-primary">
            + Nuevo Producto
        </a>
    </div>

    <div class="table-responsive">
        <table>
            <thead>
                <tr>
                    <th>ID</th>
                    <th>Producto</th>
                    <th>Stock Actual</th>
                    <th>Stock Minimo</th>
                    <th>Precio (Bs.)</th>
                    <th>Estado Stock</th>
                    <th>Acciones</th>
                </tr>
            </thead>
            <tbody>
                @forelse($inventarios as $inv)
                <tr>
                    <td>#{{ $inv->id }}</td>
                    <td><strong>{{ $inv->producto }}</strong></td>
                    <td>{{ number_format($inv->stock_actual, 2) }}</td>
                    <td>{{ number_format($inv->stock_minimo, 2) }}</td>
                    <td>Bs. {{ number_format($inv->precio, 2) }}</td>
                    <td>
                        @if($inv->stock_actual <= $inv->stock_minimo)
                            <span class="badge" style="background: rgba(239, 68, 68, 0.2); color: #f87171;">Stock Bajo</span>
                        @else
                            <span class="badge" style="background: rgba(16, 185, 129, 0.2); color: #34d399;">Normal</span>
                        @endif
                    </td>
                    <td>
                        <div class="actions-cell">
                            <a href="{{ route('inventarios.show', $inv->id) }}" class="btn btn-outline" style="padding: 0.3rem 0.6rem; font-size: 0.8rem;">Ver</a>
                            <a href="{{ route('inventarios.edit', $inv->id) }}" class="btn btn-secondary" style="padding: 0.3rem 0.6rem; font-size: 0.8rem;">Editar</a>
                            <form action="{{ route('inventarios.destroy', $inv->id) }}" method="POST" onsubmit="return confirm('Seguro de eliminar este producto del inventario?');" style="display:inline;">
                                @csrf
                                @method('DELETE')
                                <button type="submit" class="btn btn-danger" style="padding: 0.3rem 0.6rem; font-size: 0.8rem;">Eliminar</button>
                            </form>
                        </div>
                    </td>
                </tr>
                @empty
                <tr>
                    <td colspan="7" style="text-align: center; color: var(--text-muted); padding: 2rem;">
                        No hay productos registrados en el inventario.
                    </td>
                </tr>
                @endforelse
            </tbody>
        </table>
    </div>
</div>
@endsection
\end{lstlisting}

\subsection{Vista del CRUD de Inventarios en Funcionamiento}
A continuacion se presenta la captura de pantalla de la interfaz del CRUD de Inventarios:

\begin{figure}[H]
    \centering
    \includegraphics[width=0.9\textwidth]{view-inventarios.png}
    \caption{Interfaz de usuario del CRUD de Inventarios}
\end{figure}

\newpage

\section*{CRUD PEDIDOS}
\addcontentsline{toc}{section}{CRUD PEDIDOS}
\setcounter{section}{4}

\subsection{Migración (2026\_07\_22\_000003\_create\_pedidos\_table.php)}
\begin{lstlisting}[language=purePHP]
<?php

use Illuminate\Database\Migrations\Migration;
use Illuminate\Database\Schema\Blueprint;
use Illuminate\Support\Facades\Schema;

return new class extends Migration
{
    public function up(): void
    {
        Schema::create('pedidos', function (Blueprint $table) {
            $table->id();
            $table->foreignId('cliente_id')->constrained('clientes')->onDelete('cascade');
            $table->dateTime('fecha');
            $table->string('estado')->default('Pendiente');
            $table->timestamps();
        });
    }

    public function down(): void
    {
        Schema::dropIfExists('pedidos');
    }
};
\end{lstlisting}

\subsection{Modelo (Pedido.php)}
\begin{lstlisting}[language=purePHP]
<?php

namespace App\Models;

use Illuminate\Database\Eloquent\Factories\HasFactory;
use Illuminate\Database\Eloquent\Model;

class Pedido extends Model
{
    use HasFactory;

    protected $table = 'pedidos';

    protected $fillable = [
        'cliente_id',
        'fecha',
        'estado',
    ];

    public function cliente()
    {
        return $this->belongsTo(Cliente::class, 'cliente_id');
    }

    public function detalles()
    {
        return $this->hasMany(DetallePedido::class, 'pedido_id');
    }

    public function ventas()
    {
        return $this->hasMany(Venta::class, 'pedido_id');
    }

    public function getTotalAttribute()
    {
        return $this->detalles->sum('subtotal');
    }
}
\end{lstlisting}

\subsection{Controlador (PedidoController.php)}
\begin{lstlisting}[language=purePHP]
<?php

namespace App\Http\Controllers;

use App\Models\Pedido;
use App\Models\Cliente;
use Illuminate\Http\Request;

class PedidoController extends Controller
{
    public function index()
    {
        $pedidos = Pedido::with(['cliente', 'detalles.inventario'])->latest()->get();
        return view('pedidos.index', compact('pedidos'));
    }

    public function create()
    {
        $clientes = Cliente::orderBy('nombre_completo')->get();
        return view('pedidos.create', compact('clientes'));
    }

    public function store(Request $request)
    {
        $validated = $request->validate([
            'cliente_id' => 'required|exists:clientes,id',
            'fecha' => 'required|date',
            'estado' => 'required|string|in:Pendiente,En preparacion,Completado,Cancelado',
        ]);

        $pedido = Pedido::create($validated);

        return redirect()->route('pedidos.index')->with('success', 'Pedido creado con exito.');
    }

    public function show(Pedido $pedido)
    {
        $pedido->load(['cliente', 'detalles.inventario', 'ventas']);
        return view('pedidos.show', compact('pedido'));
    }

    public function edit(Pedido $pedido)
    {
        $clientes = Cliente::orderBy('nombre_completo')->get();
        return view('pedidos.edit', compact('pedido', 'clientes'));
    }

    public function update(Request $request, Pedido $pedido)
    {
        $validated = $request->validate([
            'cliente_id' => 'required|exists:clientes,id',
            'fecha' => 'required|date',
            'estado' => 'required|string|in:Pendiente,En preparacion,Completado,Cancelado',
        ]);

        $pedido->update($validated);

        return redirect()->route('pedidos.index')->with('success', 'Pedido actualizado con exito.');
    }

    public function destroy(Pedido $pedido)
    {
        $pedido->delete();

        return redirect()->route('pedidos.index')->with('success', 'Pedido eliminado con exito.');
    }
}
\end{lstlisting}

\subsection{Vista Principal (pedidos/index.blade.php)}
\begin{lstlisting}[language=HTML]
@extends('layouts.app')

@section('title', 'Pedidos')

@section('content')
<div class="glass-card">
    <div class="card-header">
        <div>
            <h1 class="card-title">Gestion de Pedidos</h1>
            <p style="color: var(--text-muted); font-size: 0.9rem;">Registro y seguimiento de comandas de clientes</p>
        </div>
        <a href="{{ route('pedidos.create') }}" class="btn btn-primary">
            + Nuevo Pedido
        </a>
    </div>

    <div class="table-responsive">
        <table>
            <thead>
                <tr>
                    <th>ID Pedido</th>
                    <th>Cliente</th>
                    <th>Fecha</th>
                    <th>Estado</th>
                    <th>Total (Bs.)</th>
                    <th>Acciones</th>
                </tr>
            </thead>
            <tbody>
                @forelse($pedidos as $pedido)
                <tr>
                    <td>#{{ $pedido->id }}</td>
                    <td><strong>{{ $pedido->cliente->nombre_completo ?? 'N/A' }}</strong></td>
                    <td>{{ \Carbon\Carbon::parse($pedido->fecha)->format('d/m/Y H:i') }}</td>
                    <td>
                        @if($pedido->estado == 'Completado')
                            <span class="badge" style="background: rgba(16, 185, 129, 0.2); color: #34d399;">Completado</span>
                        @elseif($pedido->estado == 'En preparacion')
                            <span class="badge" style="background: rgba(245, 158, 11, 0.2); color: #fbbf24;">En Preparacion</span>
                        @elseif($pedido->estado == 'Cancelado')
                            <span class="badge" style="background: rgba(239, 68, 68, 0.2); color: #f87171;">Cancelado</span>
                        @else
                            <span class="badge" style="background: rgba(79, 70, 229, 0.2); color: #a5b4fc;">Pendiente</span>
                        @endif
                    </td>
                    <td><strong>Bs. {{ number_format($pedido->total, 2) }}</strong></td>
                    <td>
                        <div class="actions-cell">
                            <a href="{{ route('pedidos.show', $pedido->id) }}" class="btn btn-outline" style="padding: 0.3rem 0.6rem; font-size: 0.8rem;">Ver Detalle</a>
                            <a href="{{ route('pedidos.edit', $pedido->id) }}" class="btn btn-secondary" style="padding: 0.3rem 0.6rem; font-size: 0.8rem;">Editar</a>
                            <form action="{{ route('pedidos.destroy', $pedido->id) }}" method="POST" onsubmit="return confirm('Seguro de eliminar este pedido?');" style="display:inline;">
                                @csrf
                                @method('DELETE')
                                <button type="submit" class="btn btn-danger" style="padding: 0.3rem 0.6rem; font-size: 0.8rem;">Eliminar</button>
                            </form>
                        </div>
                    </td>
                </tr>
                @empty
                <tr>
                    <td colspan="6" style="text-align: center; color: var(--text-muted); padding: 2rem;">
                        No hay pedidos registrados.
                    </td>
                </tr>
                @endforelse
            </tbody>
        </table>
    </div>
</div>
@endsection
\end{lstlisting}

\subsection{Vista del CRUD de Pedidos en Funcionamiento}
A continuacion se presenta la captura de pantalla de la interfaz del CRUD de Pedidos:

\begin{figure}[H]
    \centering
    \includegraphics[width=0.9\textwidth]{view-pedidos.png}
    \caption{Interfaz de usuario del CRUD de Pedidos}
\end{figure}

\newpage

\section*{CRUD DETALLE PEDIDOS}
\addcontentsline{toc}{section}{CRUD DETALLE PEDIDOS}
\setcounter{section}{5}

\subsection{Migración (2026\_07\_22\_000004\_create\_detalle\_pedidos\_table.php)}
\begin{lstlisting}[language=purePHP]
<?php

use Illuminate\Database\Migrations\Migration;
use Illuminate\Database\Schema\Blueprint;
use Illuminate\Support\Facades\Schema;

return new class extends Migration
{
    public function up(): void
    {
        Schema::create('detalle_pedidos', function (Blueprint $table) {
            $table->id();
            $table->foreignId('pedido_id')->constrained('pedidos')->onDelete('cascade');
            $table->foreignId('inventario_id')->constrained('inventarios')->onDelete('cascade');
            $table->decimal('cantidad', 10, 2);
            $table->decimal('subtotal', 10, 2);
            $table->timestamps();
        });
    }

    public function down(): void
    {
        Schema::dropIfExists('detalle_pedidos');
    }
};
\end{lstlisting}

\subsection{Modelo (DetallePedido.php)}
\begin{lstlisting}[language=purePHP]
<?php

namespace App\Models;

use Illuminate\Database\Eloquent\Factories\HasFactory;
use Illuminate\Database\Eloquent\Model;

class DetallePedido extends Model
{
    use HasFactory;

    protected $table = 'detalle_pedidos';

    protected $fillable = [
        'pedido_id',
        'inventario_id',
        'cantidad',
        'subtotal',
    ];

    public function pedido()
    {
        return $this->belongsTo(Pedido::class, 'pedido_id');
    }

    public function inventario()
    {
        return $this->belongsTo(Inventario::class, 'inventario_id');
    }
}
\end{lstlisting}

\subsection{Controlador (DetallePedidoController.php)}
\begin{lstlisting}[language=purePHP]
<?php

namespace App\Http\Controllers;

use App\Models\DetallePedido;
use App\Models\Pedido;
use App\Models\Inventario;
use Illuminate\Http\Request;

class DetallePedidoController extends Controller
{
    public function index()
    {
        $detalles = DetallePedido::with(['pedido.cliente', 'inventario'])->latest()->get();
        return view('detalle_pedidos.index', compact('detalles'));
    }

    public function create(Request $request)
    {
        $pedidos = Pedido::with('cliente')->latest()->get();
        $inventarios = Inventario::where('stock_actual', '>', 0)->get();
        $selected_pedido_id = $request->query('pedido_id');

        return view('detalle_pedidos.create', compact('pedidos', 'inventarios', 'selected_pedido_id'));
    }

    public function store(Request $request)
    {
        $validated = $request->validate([
            'pedido_id' => 'required|exists:pedidos,id',
            'inventario_id' => 'required|exists:inventarios,id',
            'cantidad' => 'required|numeric|min:0.01',
            'subtotal' => 'nullable|numeric|min:0',
        ]);

        $inventario = Inventario::findOrFail($validated['inventario_id']);

        if (!isset($validated['subtotal']) || $validated['subtotal'] == 0) {
            $validated['subtotal'] = $validated['cantidad'] * $inventario->precio;
        }

        $detalle = DetallePedido::create($validated);
        $inventario->decrement('stock_actual', $validated['cantidad']);

        return redirect()->route('pedidos.show', $validated['pedido_id'])
            ->with('success', 'Detalle de pedido agregado con exito.');
    }

    public function show(DetallePedido $detallePedido)
    {
        $detallePedido->load(['pedido.cliente', 'inventario']);
        return view('detalle_pedidos.show', compact('detallePedido'));
    }

    public function edit(DetallePedido $detallePedido)
    {
        $pedidos = Pedido::with('cliente')->get();
        $inventarios = Inventario::all();
        return view('detalle_pedidos.edit', compact('detallePedido', 'pedidos', 'inventarios'));
    }

    public function update(Request $request, DetallePedido $detallePedido)
    {
        $validated = $request->validate([
            'pedido_id' => 'required|exists:pedidos,id',
            'inventario_id' => 'required|exists:inventarios,id',
            'cantidad' => 'required|numeric|min:0.01',
            'subtotal' => 'nullable|numeric|min:0',
        ]);

        $inventario = Inventario::findOrFail($validated['inventario_id']);

        if (!isset($validated['subtotal']) || $validated['subtotal'] == 0) {
            $validated['subtotal'] = $validated['cantidad'] * $inventario->precio;
        }

        $detallePedido->update($validated);

        return redirect()->route('pedidos.show', $validated['pedido_id'])
            ->with('success', 'Detalle de pedido actualizado con exito.');
    }

    public function destroy(DetallePedido $detallePedido)
    {
        $pedidoId = $detallePedido->pedido_id;
        $detallePedido->delete();

        return redirect()->route('pedidos.show', $pedidoId)
            ->with('success', 'Detalle de pedido eliminado con exito.');
    }
}
\end{lstlisting}

\subsection{Vista del CRUD de Detalle Pedidos en Funcionamiento}
A continuacion se presenta la captura de pantalla de la interfaz del CRUD de Detalle Pedidos:

\begin{figure}[H]
    \centering
    \includegraphics[width=0.9\textwidth]{view-detalle-pedidos.png}
    \caption{Interfaz de usuario del CRUD de Detalle Pedidos}
\end{figure}

\newpage

\section*{CRUD VENTAS}
\addcontentsline{toc}{section}{CRUD VENTAS}
\setcounter{section}{6}

\subsection{Migración (2026\_07\_22\_000005\_create\_ventas\_table.php)}
\begin{lstlisting}[language=purePHP]
<?php

use Illuminate\Database\Migrations\Migration;
use Illuminate\Database\Schema\Blueprint;
use Illuminate\Support\Facades\Schema;

return new class extends Migration
{
    public function up(): void
    {
        Schema::create('ventas', function (Blueprint $table) {
            $table->id();
            $table->foreignId('pedido_id')->constrained('pedidos')->onDelete('cascade');
            $table->dateTime('fecha');
            $table->decimal('total', 10, 2);
            $table->timestamps();
        });
    }

    public function down(): void
    {
        Schema::dropIfExists('ventas');
    }
};
\end{lstlisting}

\subsection{Modelo (Venta.php)}
\begin{lstlisting}[language=purePHP]
<?php

namespace App\Models;

use Illuminate\Database\Eloquent\Factories\HasFactory;
use Illuminate\Database\Eloquent\Model;

class Venta extends Model
{
    use HasFactory;

    protected $table = 'ventas';

    protected $fillable = [
        'pedido_id',
        'fecha',
        'total',
    ];

    public function pedido()
    {
        return $this->belongsTo(Pedido::class, 'pedido_id');
    }
}
\end{lstlisting}

\subsection{Controlador (VentaController.php)}
\begin{lstlisting}[language=purePHP]
<?php

namespace App\Http\Controllers;

use App\Models\Venta;
use App\Models\Pedido;
use Illuminate\Http\Request;

class VentaController extends Controller
{
    public function index()
    {
        $ventas = Venta::with(['pedido.cliente', 'pedido.detalles.inventario'])->latest()->get();
        return view('ventas.index', compact('ventas'));
    }

    public function create()
    {
        $pedidos = Pedido::with('cliente')->latest()->get();
        return view('ventas.create', compact('pedidos'));
    }

    public function store(Request $request)
    {
        $validated = $request->validate([
            'pedido_id' => 'required|exists:pedidos,id',
            'fecha' => 'required|date',
            'total' => 'required|numeric|min:0',
        ]);

        $venta = Venta::create($validated);

        return redirect()->route('ventas.index')->with('success', 'Venta registrada con exito.');
    }

    public function show(Venta $venta)
    {
        $venta->load(['pedido.cliente', 'pedido.detalles.inventario']);
        return view('ventas.show', compact('venta'));
    }

    public function edit(Venta $venta)
    {
        $pedidos = Pedido::with('cliente')->latest()->get();
        return view('ventas.edit', compact('venta', 'pedidos'));
    }

    public function update(Request $request, Venta $venta)
    {
        $validated = $request->validate([
            'pedido_id' => 'required|exists:pedidos,id',
            'fecha' => 'required|date',
            'total' => 'required|numeric|min:0',
        ]);

        $venta->update($validated);

        return redirect()->route('ventas.index')->with('success', 'Venta actualizada con exito.');
    }

    public function destroy(Venta $venta)
    {
        $venta->delete();

        return redirect()->route('ventas.index')->with('success', 'Venta eliminada con exito.');
    }
}
\end{lstlisting}

\subsection{Vista Principal (ventas/index.blade.php)}
\begin{lstlisting}[language=HTML]
@extends('layouts.app')

@section('title', 'Ventas')

@section('content')
<div class="glass-card">
    <div class="card-header">
        <div>
            <h1 class="card-title">Registro de Ventas</h1>
            <p style="color: var(--text-muted); font-size: 0.9rem;">Facturacion y cierre de ventas realizadas</p>
        </div>
        <a href="{{ route('ventas.create') }}" class="btn btn-primary">
            + Registrar Venta
        </a>
    </div>

    <div class="table-responsive">
        <table>
            <thead>
                <tr>
                    <th>ID Venta</th>
                    <th>Pedido ID</th>
                    <th>Cliente</th>
                    <th>Fecha</th>
                    <th>Total (Bs.)</th>
                    <th>Acciones</th>
                </tr>
            </thead>
            <tbody>
                @forelse($ventas as $venta)
                <tr>
                    <td>#{{ $venta->id }}</td>
                    <td>
                        <a href="{{ route('pedidos.show', $venta->pedido_id) }}" style="color: var(--primary-color); text-decoration: underline;">
                            Pedido #{{ $venta->pedido_id }}
                        </a>
                    </td>
                    <td><strong>{{ $venta->pedido->cliente->nombre_completo ?? 'N/A' }}</strong></td>
                    <td>{{ \Carbon\Carbon::parse($venta->fecha)->format('d/m/Y H:i') }}</td>
                    <td><strong>Bs. {{ number_format($venta->total, 2) }}</strong></td>
                    <td>
                        <div class="actions-cell">
                            <a href="{{ route('ventas.show', $venta->id) }}" class="btn btn-outline" style="padding: 0.3rem 0.6rem; font-size: 0.8rem;">Ver</a>
                            <a href="{{ route('ventas.edit', $venta->id) }}" class="btn btn-secondary" style="padding: 0.3rem 0.6rem; font-size: 0.8rem;">Editar</a>
                            <form action="{{ route('ventas.destroy', $venta->id) }}" method="POST" onsubmit="return confirm('Seguro de eliminar esta venta?');" style="display:inline;">
                                @csrf
                                @method('DELETE')
                                <button type="submit" class="btn btn-danger" style="padding: 0.3rem 0.6rem; font-size: 0.8rem;">Eliminar</button>
                            </form>
                        </div>
                    </td>
                </tr>
                @empty
                <tr>
                    <td colspan="6" style="text-align: center; color: var(--text-muted); padding: 2rem;">
                        No hay ventas registradas.
                    </td>
                </tr>
                @endforelse
            </tbody>
        </table>
    </div>
</div>
@endsection
\end{lstlisting}

\subsection{Vista del CRUD de Ventas en Funcionamiento}
A continuacion se presenta la captura de pantalla de la interfaz del CRUD de Ventas:

\begin{figure}[H]
    \centering
    \includegraphics[width=0.9\textwidth]{view-ventas.png}
    \caption{Interfaz de usuario del CRUD de Ventas}
\end{figure}

\newpage

\section*{INFORME TÉCNICO: FUNCIONAMIENTO DE MODELOS Y CONTROLADORES (CASO EMPLEADOS)}
\addcontentsline{toc}{section}{INFORME TÉCNICO: FUNCIONAMIENTO DE MODELOS Y CONTROLADORES (CASO EMPLEADOS)}
\setcounter{section}{7}

\subsection{Introducción al Patrón MVC en Laravel}
El framework Laravel implementa el patrón de arquitectura software \textbf{Modelo-Vista-Controlador (MVC)}, el cual promueve una clara separación de responsabilidades:

\begin{itemize}
    \item \textbf{Modelo (\texttt{Empleado.php}):} Representa la capa de datos y lógica de negocio. Mapea la tabla \texttt{empleados} en MySQL mediante Eloquent ORM.
    \item \textbf{Controlador (\texttt{EmpleadoController.php}):} Actúa como mediador. Procesa las peticiones HTTP del usuario, aplica validaciones de seguridad y coordina las consultas con el modelo para retornar la respuesta visual.
    \item \textbf{Vista (\texttt{resources/views/empleados/}):} Renderiza la interfaz gráfica en HTML/Blade presentada al cliente.
\end{itemize}

\subsection{Funcionamiento Interno del Modelo (\texttt{Empleado.php})}
El modelo \texttt{Empleado} hereda de la clase base \texttt{Illuminate\textbackslash Database\textbackslash Eloquent\textbackslash Model}.

\subsubsection{Asignación Masiva (\texttt{\$fillable})}
La propiedad protegida \texttt{\$fillable} especifica explícitamente qué campos del formulario pueden ser insertados o actualizados mediante métodos como \texttt{Empleado::create(\$data)}. Esto previene ataques de \textit{Mass Assignment Violation}.

\subsubsection{Relaciones entre Entidades (\texttt{belongsTo})}
La función \texttt{turno()} define una relación de pertenencia hacia el modelo \texttt{Turno}:
\begin{lstlisting}[language=purePHP]
public function turno()
{
    return $this->belongsTo(Turno::class, 'id_turno');
}
\end{lstlisting}
Esta relación permite acceder directamente a los atributos del turno desde la instancia de un empleado (por ejemplo: \texttt{\$empleado->turno->dias\_descanso}).

\subsection{Funcionamiento del Controlador (\texttt{EmpleadoController.php})}
El controlador de recursos agrupa las 7 acciones estándar de un CRUD RESTful:

\begin{enumerate}
    \item \textbf{\texttt{index()}:} Recupera la lista de empleados de la base de datos aplicando \textit{Eager Loading} (\texttt{Empleado::with('turno')->get()}) para optimizar las consultas SQL eliminando el problema $N+1$, y retorna la vista de listado.
    \item \textbf{\texttt{create()}:} Consulta todos los turnos disponibles (\texttt{Turno::all()}) y renderiza el formulario para registrar un nuevo empleado.
    \item \textbf{\texttt{store(Request \$request)}:} Ejecuta la validación estricta de la solicitud HTTP (\texttt{\$request->validate()}), verificando la unicidad del CI y correo electrónico, y la existencia previa del \texttt{id\_turno}. Si es correcto, persiste el nuevo empleado en MySQL.
    \item \textbf{\texttt{show(Empleado \$empleado)}:} Utiliza \textit{Route Model Binding} para inyectar automáticamente la instancia del empleado correspondiente al ID recibido en la URL y mostrar sus detalles.
    \item \textbf{\texttt{edit(Empleado \$empleado)} y \texttt{update(Request \$request, Empleado \$empleado)}:} Muestran y procesan la edición del empleado, ignorando el propio ID del usuario durante la validación de unicidad (\texttt{unique:empleados,ci,' . \$empleado->id}).
    \item \textbf{\texttt{destroy(Empleado \$empleado)}:} Elimina el registro del empleado de manera segura y emite una notificación de éxito vía sesión.
\end{enumerate}

\subsection{Flujo Paso a Paso de una Petición de Empleados}
\begin{enumerate}
    \item El usuario ingresa a la URL \texttt{/empleados} en el navegador.
    \item El enrutador (\texttt{routes/web.php}) intercepta la petición y llama al método \texttt{EmpleadoController@index}.
    \item El controlador solicita los datos al modelo mediante \texttt{Empleado::with('turno')->get()}.
    \item Eloquent ejecuta la consulta SQL hacia la base de datos MySQL y retorna los objetos estructurados.
    \item El controlador transmite la variable \texttt{\$empleados} a la vista Blade \texttt{empleados.index}.
    \item Laravel compila la plantilla Blade y responde con el HTML estilizado con Tailwind CSS al usuario.
\end{enumerate}

\newpage
\section*{REPOSITORIO DEL PROYECTO}
\addcontentsline{toc}{section}{REPOSITORIO DEL PROYECTO}

Todo el código fuente desarrollado para esta entrega se encuentra actualizado y disponible en el repositorio de GitHub:

\begin{center}
    \url{https://github.com/Chanka-Dev/pollos-luisa}
\end{center}

\end{document}
